% Final project report (deliverable D14). Build (from this folder):
%   python -m src.viz.static_maps && python -m src.viz.charts   (from the repo root, once)
%   pdflatex report && bibtex report && pdflatex report && pdflatex report
% Figures are read from outputs/figures/ (git-ignored, regenerated by src/viz).
\documentclass[11pt,a4paper]{article}

\usepackage[margin=2.4cm]{geometry}
\usepackage[T1]{fontenc}
\usepackage{lmodern}
\usepackage{textcomp}
\usepackage{microtype}
\usepackage{graphicx}
\usepackage{booktabs}
\usepackage{tabularx}
\usepackage{amsmath}
\usepackage{caption}
\usepackage{subcaption}
\usepackage[round,authoryear]{natbib}
\usepackage{xcolor}
\usepackage[colorlinks=true,linkcolor=black,citecolor=black,urlcolor=blue]{hyperref}

\graphicspath{{../../outputs/figures/}}
% units without siunitx (not installed everywhere): \SI{3}{\km}, \SIrange{1}{2}{\m}, \ang{15}
\newcommand{\km}{km}
\newcommand{\m}{m}
\newcommand{\hectare}{ha}
\newcommand{\squared}{\textsuperscript{2}}
\newcommand{\SI}[2]{#1\,#2}
\newcommand{\SIrange}[3]{#1--#2\,#3}
\newcommand{\ang}[1]{#1\textdegree}
\captionsetup{font=small,labelfont=bf}
\newcommand{\todo}[1]{\textcolor{red!70!black}{\textit{[To be added: #1]}}}

\title{GIS-Based Land Suitability Analysis for Urban Development\\
       using Land Use Land Cover: Bengaluru South}
\author{Harsh Chattar (22CS30028) \and Abhinav (22CS30005)\\[0.3em]
  \small CS60111 Geographical Information Systems}
\date{October 2026}

\begin{document}
\maketitle

\begin{abstract}
Where can new urban development reasonably go? We answer this question for a
582\,km\textsuperscript{2} area of south Bengaluru, India, entirely in Python, from open
10\,m land-cover maps (2018--2023), a 30\,m DEM and OpenStreetMap. Land is classified into
built-up, forest and usable by K-Means clustering of 100\,m cells (96.6\,\% agreement with the
source classes). For any existing building, the usable cells whose surroundings (land cover in
250 and 500\,m rings, distances to built-up land, roads and water, terrain) are most similar
are found and explained. A building-similarity score and a random forest trained on recent
change are combined into a suitability score, graded into classes and grouped into candidate
sites. Every score is computed with data as of 2018 and tested against land that became built
in both 2022 and 2023, under strict leakage rules and next to simple baselines. The combined score
reaches an AUC of 0.797 (top-decile lift 3.45), ahead of the distance to built-up land (0.763)
and a classic AHP overlay (0.636); its two best classes hold 65\,\% of the growth. Growth next
to existing built-up land is predicted well, growth away from it not at all.
\end{abstract}

\tableofcontents
\newpage

\input{sections/01_introduction}
% Owner: Abhinav (P8.5); synthesis written with Harsh (docs/literature_review.md)
\section{Literature review}\label{sec:literature}

Two families of methods rank land for urban development: expert MCDA and learned suitability.
Table~\ref{tab:literature} summarises the papers reviewed; the full notes are in
\texttt{docs/literature\_review.md}.

\subsection{Expert multi-criteria suitability}

Land-use suitability analysis began with hand-drawn overlays and moved to GIS, while its core
idea stayed the same: map the factors, rate them, and overlay them into a composite
\citep{collins2001}. \citet{malczewski2004} groups GIS methods into overlay mapping, MCDA
(weighted linear combination, often with AHP weights) and soft-computing approaches. He notes
that results depend heavily on subjective choices of criteria, standardisation and weights, and
that suitability maps are rarely validated. The weighting method alone matters: on the same
criteria, AHP and Fuzzy-AHP ranked development options the same way but drew different zone
boundaries \citep{mosadeghi2015}. A typical Indian study, \citet{parry2018}, combines slope,
altitude, LULC and the status of existing amenities with AHP weights at ward level for Srinagar
and Jammu, without validation.

The criteria are similar across studies: terrain (slope, elevation), current land cover, and
proximity to roads, existing built-up land and water. Our feature table has the same core, at
100\,m cells instead of wards, and our MCDA baseline uses exactly these criteria.

\subsection{LULC change in Bengaluru and Indian cities}

\citet{ramachandra2012} classified Landsat images of Bangalore from 1973 to 2010: built-up area
grew 584\,\% while vegetation fell 66\,\% and water bodies 74\,\%. Growth has eaten into exactly
the land our exclusion mask protects (lakes and forest). Across five other Indian mega cities,
\citet{bharath2018} find that industrial, infrastructural and socio-economic drivers influence
growth more than biophysical factors. This matches our design, in which road access and nearby
built-up land are the main context features and terrain is secondary.

\subsection{Learning suitability from data}

A second line of work learns a suitability surface from where development already is or recently
happened: neural networks \citep{pijanowski2002, liu2017flus}, random forests
\citep{ahmadlou2016, li2024} and presence-only maximum entropy \citep{wang2021}. They share a
core of predictors with ours: distances to roads, rivers and existing urban land, terrain, and
neighbourhood land use.

\textbf{Validation is the weak point.} Studies that test on a random split of a single date
report very high accuracy, such as AUC 0.99 for urban suitability in Yulin \citep{li2024}, where
the cell's own land use is also an input. Studies that test against \emph{later observed change}
report more modest values, such as AUC 0.82 in Rasht \citep{ahmadlou2016}. Random
cross-validation over-states performance on spatially structured data
\citep{roberts2017, ploton2020}; spatial blocks give more honest estimates \citep{valavi2019}.
Growth type also matters: growth away from existing urban land is much harder to predict than
growth next to it \citep{wang2021}, which can be measured with the Landscape Expansion Index
\citep{liu2010lei}. The total operating characteristic complements ROC for maps where change is
a small share of the land \citep{pontius2014}.

\textbf{Data quality} shapes the design as well. Among the global 10\,m LULC products, the Esri
map has the highest overall accuracy, but it over-estimates scrub \citep{venter2022}. OSM
building completeness varies strongly between cities and over time \citep{herfort2023}.

\subsection{The gap this project addresses}

Expert MCDA is transparent but subjective and almost never tested against real growth. Learned
models remove the hand-set weights but are often validated optimistically. This project:
\begin{itemize}
  \item tests \emph{every} score, including the expert MCDA baseline, the same way: against
        persistent growth from 2018/19 to 2022/23, always next to simple baselines (random and
        distance to built-up);
  \item enforces leakage rules: no own-cell land cover in model inputs, OSM data as of the
        baseline date, a time-travel rule, and out-of-fold scores with spatial blocks;
  \item adds a per-building similarity query (``find land whose surroundings resemble
        building~A''), which none of the reviewed papers offers.
\end{itemize}

\begin{table}[ht]
  \centering\small
  \caption{Papers reviewed. Validation: how the paper tests its result.}
  \label{tab:literature}
  \begin{tabularx}{\textwidth}{@{}L{3.3cm} L{2.3cm} Y Y@{}}
    \toprule
    Paper & Area & Method & Validation \\
    \midrule
    \citet{malczewski2004} & review & GIS-MCDA methods & rarely done (noted) \\
    \citet{collins2001} & USA (review) & history of suitability analysis & --- \\
    \citet{mosadeghi2015} & Gold Coast & AHP vs Fuzzy-AHP overlay & two outputs compared \\
    \citet{parry2018} & Srinagar, Jammu & GIS-AHP, ward level & none reported \\
    \citet{ramachandra2012} & Bangalore & Landsat LULC change, 1973--2010 & classification accuracy \\
    \citet{bharath2018} & 5 Indian cities & CA-Markov with growth agents & not reported in the abstract \\
    \citet{li2024} & Yulin & random forest, 25 factors & random 70/30 split, AUC 0.99 \\
    \citet{wang2021} & Luoyang & MaxEnt + cellular automaton & 2009$\to$2018, Kappa 0.78 \\
    \citet{ahmadlou2016} & Rasht & random forest & later change, AUC 0.82, TOC \\
    \citet{liu2017flus} & China & neural network + CA (FLUS) & 2000$\to$2010 simulation \\
    \citet{pijanowski2002} & Michigan & neural network (LTM), 100\,m & same period \\
    \citet{roberts2017}, \citet{ploton2020} & methods & blocked cross-validation & random vs blocked CV \\
    \citet{liu2010lei}, \citet{pontius2014} & methods & LEI growth types; TOC curve & --- \\
    \citet{venter2022}, \citet{herfort2023} & global & LULC accuracy; OSM completeness & ground truth; ML estimate \\
    \bottomrule
  \end{tabularx}
\end{table}

\input{sections/03_study_area}
% Owner: Abhinav (P8.5); details in docs/data_sources.md
\section{Data}\label{sec:data}

All data are free and openly licensed (Table~\ref{tab:data}). Rasters are read from the Microsoft
Planetary Computer, and only the window needed for the study area plus a \SI{3}{\km} buffer is
downloaded. The buffer equals the distance cap of the features, so distances from cells near the
edge are exact. OpenStreetMap layers come from the Overpass API through OSMnx \citep{boeing2017}.
Every download is recorded with its checksum and extent in a manifest, so a rerun reproduces or
skips it.

\begin{table}[ht]
  \centering\small
  \caption{Datasets.}
  \label{tab:data}
  \begin{tabularx}{\textwidth}{@{}Y L{1.9cm} L{1.6cm} L{2.4cm} Y@{}}
    \toprule
    Dataset & Years & Resolution & Licence & Used for \\
    \midrule
    ESRI 10\,m Annual LULC v02 \citep{karra2021} & 2018--2023 & \SI{10}{\m} & CC BY 4.0 &
      land-cover features, three-class map, growth labels \\
    ESA WorldCover v200 \citep{zanaga2022} & 2021 & \SI{10}{\m} & CC BY 4.0 &
      cross-check of the three-class map \\
    Copernicus DEM GLO-30 \citep{copdem} & static & \SI{30}{\m} & Copernicus DEM licence &
      elevation, slope \\
    OpenStreetMap \citep{osm} & 2018-01-01, current & vector & ODbL 1.0 &
      roads, water, buildings, protected areas \\
    \bottomrule
  \end{tabularx}
\end{table}

\paragraph{Land use / land cover.} The ESRI (Impact Observatory) map is used for every year and
every purpose, so that ``built'' has the same definition in the baseline and the test years. It
has the highest overall accuracy of the global 10\,m products (75\,\%), and built area is one of
its most reliable classes \citep{venter2022}. Its known bias, over-estimating scrub, matters
here: much of the Bannerghatta forest appears as rangeland. Inside the park, ESRI classes 47\,\%
of the land as trees and 50\,\% as rangeland. ESA WorldCover 2021 is used only as an independent
cross-check. The two products disagree strongly on built-up land (62\,\% of the area in ESRI 2023
vs 33\,\% in WorldCover), because ESRI generalises urban areas to include gardens and street
trees. ESRI land cover also flickers between years: only 35--49\,\% of single-year ``new built''
pixels are still built later. The growth labels therefore require two years at each end
(Section~\ref{sec:validation_design}).

\paragraph{Terrain.} The Copernicus DEM is warped to UTM at \SI{30}{\m}. Slope is computed with
Horn's method \citep{horn1981} and then resampled to the \SI{10}{\m} grid. The DEM is a surface
model, so heights include buildings and trees; slope is less reliable in dense built-up cells, but
it matters most on open candidate land, where this effect is small.

\paragraph{OpenStreetMap.} Roads, water and buildings are downloaded twice: as they were on
1 January 2018 (an Overpass \texttt{[date:]} query) for the validation run, and as they are now
for the final map. The difference is large. All-road length more than doubled between the two
snapshots (\SI{3889}{\km} to \SI{8730}{\km}), partly new streets in growth areas and partly later
mapping. Road features from today's map would therefore leak future growth into a 2018 analysis.
OSM building coverage is also very uneven: buildings were mapped in 62--75\,\% of the built cells
in the north but in about 1\,\% in the south in 2018, and almost all are untagged
(\texttt{building=yes}). Protected areas (Bannerghatta, tagged as a national park) are a static
layer and come from the current map.

\paragraph{Preprocessing.} All rasters are aligned to one reference grid in EPSG:32643: \SI{10}{\m}
pixels whose edges fall on whole \SI{100}{\m} cells, so every block of $10\times10$ pixels is one
grid cell. Land cover is copied without resampling, because the ESRI tile is already on this
lattice. OSM layers are reprojected, clipped to the grid, repaired and split into single-part
geometries. Drivable roads are split into major (motorway to secondary) and minor. Buildings
smaller than \SI{10}{\m\squared} or larger than \SI{10}{\hectare} are dropped. Each preparation
step is checked independently against the source data, for example by comparing pixel values at
random points, matching road lengths before and after clipping, and overlaying OSM lakes on ESRI
water.

% Methodology (P8.6)
\section{Methodology}\label{sec:methodology}

The analysis runs as one pipeline (\texttt{python -m src.pipeline}). Every stage reads and
writes files with a fixed schema, and every parameter lives in \texttt{config/config.yaml}.

\subsection{Analysis grid and features}

The study area plus a \SI{3}{\km} buffer is divided into \SI{100}{\m} cells aligned with the
\SI{10}{\m} raster grid, so every cell holds exactly $10\times10$ land-cover pixels. The
58{,}218 cells whose centres lie inside the study area are analysed. For each cell and year
the feature table holds:
\begin{itemize}
  \item \textbf{own-cell land cover}: the shares of water, trees, flooded vegetation, crops,
        built area, bare ground and rangeland among its 100 pixels;
  \item \textbf{context rings}: the same shares in square rings of \SI{250}{\m} and
        \SI{500}{\m} around the cell, \emph{excluding} the cell itself;
  \item \textbf{distances} from the cell centre to the nearest built-up pixel outside the cell,
        to water, to a major road and to any road (log-transformed, capped at \SI{3}{\km}, the
        buffer width, so that edge cells see everything within the cap);
  \item \textbf{road density}: road length within \SI{500}{\m} per km\textsuperscript{2};
  \item \textbf{terrain}: mean elevation, mean and maximum slope;
  \item building count and footprint share (kept for analysis only).
\end{itemize}

\textbf{Leakage rule.} A model trained to recognise built land could simply read it from the
cell's own land cover, whose shares sum to one minus the built share, or from its buildings.
In a prototype, own-cell shares alone separated built from non-built cells perfectly. The
similarity and the random forests therefore use only the context features: rings, distances,
road density and terrain (20 inputs). The code refuses any other column.

\subsection{Land classification}

K-Means \citep{pedregosa2011} groups the cells on seven standardised inputs: the shares of water,
trees, crops, built area and rangeland, mean elevation and mean slope. (Snow is absent and
flooded vegetation and bare ground are almost always zero; maximum slope correlates at
$r=0.95$ with mean slope.) The number of clusters is chosen from $k=3$ to 10 with inertia, the
silhouette \citep{rousseeuw1987} and the Davies--Bouldin index \citep{davies1979}; Gaussian
mixtures are compared. Clusters become classes by rules on their centroids:
water-dominated $\rightarrow$ excluded; built-up if built land exceeds 30\,\% and is the largest
land group; forest if trees exceed 50\,\%, or trees and rangeland together exceed 50\,\% on
slopes above \ang{8} (the Bannerghatta scrub forest, which the ESRI map partly calls rangeland);
otherwise usable, with cropland flagged as agricultural.

The \textbf{exclusion mask} is applied last, cell by cell: more than half of the cell is water,
flooded vegetation or snow in either year of the run, or an OSM lake; mean slope above
\ang{15}; more than half no-data; or more than half inside a protected area.

\subsection{Building similarity}\label{sec:similarity}

\paragraph{Reference buildings.} Cells with an OpenStreetMap building, plus ESRI built cells
where OSM has none: OSM maps 62--75\,\% of built cells in the north but about 1\,\% in the south
in 2018, so OSM alone would describe only northern surroundings.

\paragraph{Profiles.} A location's profile is its cell's 20 context features, standardised with
the mean and standard deviation of the reference and candidate cells of the run.

\paragraph{Single query.} For building A, every usable cell is ranked by the Euclidean distance
between its profile and A's (cosine distance is available for comparison). Each result lists
the features where B is closest to A and the ones where it differs most, in standard deviations.

\paragraph{Aggregate similarity.} For every candidate cell, the mean distance to its ten nearest
reference profiles, never counting the cell itself, is rescaled to 0--1 between the 1st and
99th percentile (higher = more similar).

\subsection{Learned models}

Two random forests \citep{breiman2001} use the 20 context features of 2018:
\begin{itemize}
  \item the \textbf{change RF} learns from the cells that were non-built in 2018 and 2019 and
        became built in 2020 and 2021 (355 positives among 15{,}658 candidates);
  \item the \textbf{state RF} learns built vs non-built cells in 2018 (for comparison).
\end{itemize}
Both are class-balanced (300 trees). Hyperparameters are tuned by spatial cross-validation
\citep{roberts2017, valavi2019}: cells are grouped into \SI{2}{\km} blocks, which are dealt into
five folds, so no block is in both training and test data. Tuning uses only the training label.
Features enter as percentiles within their year, because OSM road mapping grew a lot between
the 2018 and current snapshots; within a year the transform keeps each tree's partition of its
training data. Importance is reported both by impurity and by permutation on the held-out folds.

\subsection{MCDA / AHP baseline}

The classic overlay rescales five criteria to 0--1: distance to built-up land (1 next to it, 0
beyond \SI{2}{\km}), distance to a major road (0 beyond \SI{3}{\km}), slope (0 at \ang{15}),
own-cell land cover (rangeland and bare ground 1, crops 0.6, trees 0.2) and a water setback (0
within \SI{100}{\m}, 1 beyond \SI{200}{\m}). Weights come from an AHP pairwise matrix
\citep{saaty1980}: built-up and road distance 0.30 each, slope and land cover 0.16 each, water
0.09, with a consistency ratio of 0.003. MCDA is not trained, so its use of own-cell land cover
is not the leakage problem above.

\subsection{Combined score, classes and candidate sites}

The final score is a weighted mean of the aggregate similarity and the change RF. Each is first
converted to its percentile among the scored cells, because the two have very different
distributions. The weight is chosen on the validation run, but on the change RF's training label
(growth by 2020/21), never on the later growth used for validation, and then frozen: 0.3 for the
similarity and 0.7 for the change RF.

Usable cells are classed by quantiles of the score: S1 at least the 90th percentile, S2 the 70th,
S3 the 40th, otherwise N. Forest, built-up and excluded land is always N. Edge-adjacent S1 cells
are merged into candidate patches of at least \SI{1}{\hectare}, each with its area, mean score,
slope, distances and land cover. The top sites are the 20 best patches of at least
\SI{5}{\hectare}.

\subsection{Implementation}

Everything is written in Python: rasterio and GeoPandas for the data, scikit-learn for the
models, Matplotlib and Folium for the maps; OpenStreetMap data comes through OSMnx
\citep{boeing2017}. Each stage has unit tests on synthetic data, a phase gate on the real data,
and an independent check script that recomputes its output another way (for example pixel
values against the source rasters, brute-force distances, or AUC from ranks). The whole pipeline
runs in about 12 minutes after the first download, and a fresh clone reproduces the validation
run exactly.

\paragraph{Live demo.} The results are also published as a static website:
\begin{center}
  \url{https://iamabhinav-dev.github.io/gis-urban-suitability/}
\end{center}
It shows the suitability map with the top 20 sites, an interactive map of every layer, and a
``find places like A'' tool. The tool runs the single-building query of Section~\ref{sec:similarity}
in the browser on the exported cell profiles; a check script confirms that it returns the same
top 20 locations as the Python query.

% Validation design (P8.6)
\section{Validation design}\label{sec:validation_design}

The question is not whether a score reproduces today's land cover, which is easy, but whether
it would have pointed to the land that was built on next. Every score is therefore computed
twice.

\paragraph{Two runs.} The \textbf{validation run} uses only data available in 2019 or earlier:
ESRI land cover of 2018, OpenStreetMap as of 1 January 2018, reference buildings of 2018, and a
three-class map and exclusion mask built from 2018 and 2019. Its scores are tested against the
growth that followed. The \textbf{final run} applies the same models to the latest data (2023
land cover, OSM as of 1 October 2026) to produce the suitability map. Two layers are treated as
static and used in both runs: the DEM, and the boundary of the national park.

\paragraph{Growth labels.} ESRI land cover flickers between years: only 35--49\,\% of the pixels
that turn built in one year stay built. Labels therefore require two years at each end. A
\textbf{candidate} cell has a built share below 10\,\% in both 2018 and 2019 and is not
excluded. It \textbf{grew} if its built share is at least 50\,\% in both 2022 and 2023.
Candidates that end up in between, or flicker, are left out (2{,}881 cells). This gives
12{,}777 candidate cells, of which 1{,}125 grew (8.8\,\%). Each grown cell is also typed by the
landscape expansion index \citep{liu2010lei}: \emph{adjacent} growth touches built-up land
that existed in 2018, \emph{outlying} growth does not.

\paragraph{Leakage rules.}
\begin{itemize}
  \item No model sees a cell's own land cover or buildings (Section~\ref{sec:methodology}).
  \item Nothing dated after 2019 feeds a validation-run score.
  \item The change RF learns from growth by 2020/21, and most of those cells also count as
        growth by 2022/23. Its validation scores are therefore \textbf{out of fold}: each cell
        is scored by a model trained without its \SI{2}{\km} block. In-sample scores would
        reach AUC 0.997.
  \item The combined score's weight is chosen on the 2020/21 label, not on the test label.
\end{itemize}

\paragraph{Metrics.} One harness scores every model the same way:
\begin{itemize}
  \item ROC-AUC on the 12{,}777 candidate cells;
  \item precision, recall and \textbf{lift} (precision divided by the 8.8\,\% growth rate) in
        the top 10\,\% and 30\,\% of scores;
  \item the total operating characteristic \citep{pontius2014}, as reported by
        \citet{ahmadlou2016};
  \item AUC separately for adjacent and outlying growth;
  \item for the combined score, the share of growth that falls in classes S1 and S2.
\end{itemize}
Random scores and the distance to built-up land are always reported next to the models: in a
prototype, distance to built-up alone reached AUC 0.75, so it is the bar to beat. The targets set
in advance were a best model at least as good as this baseline and a lift of at least 2 in the
top 10\,\%.

\paragraph{Further checks.}
\begin{itemize}
  \item \textbf{Ablation:} each feature group is dropped in turn.
  \item \textbf{Held-out buildings:} reference cells hidden by spatial block must be ranked high
        by the rest.
  \item \textbf{Self-query:} querying a building must return its own location near the top.
  \item \textbf{Sensitivity:} the analysis is repeated with slope thresholds of \ang{10} and
        \ang{20}, $k=5$ to 8, and cells of 50 and \SI{200}{\m}.
\end{itemize}

% Results (P8.7)
\section{Results}\label{sec:results}

\subsection{Land classification}\label{sec:res-classification}

\paragraph{Choosing the clusters.} K-Means and Gaussian mixture models were run for $k=3$ to 10
on seven standardised inputs: the cell's shares of water, trees, crops, built and rangeland, its
mean elevation and its mean slope (Figure~\ref{fig:kscan}). K-Means separates the cells far
better at every $k$: silhouette 0.43--0.61 \citep{rousseeuw1987}, against 0.26--0.36 for the
mixture models, which also have a worse Davies--Bouldin index \citep{davies1979} at every $k$. The
silhouette of K-Means is flat from $k=5$ to $k=8$ and drops sharply at $k=9$. We use
\textbf{$\boldsymbol{k=7}$} (silhouette 0.606 on the 2023 features, 0.585 on 2018), which is the
best or within 0.003 of the best in both years. The seven clusters are the five land-cover groups
plus two sub-groups that appear in both years:
\begin{itemize}
  \item \textbf{steep natural vegetation} (mean slope \ang{13.2}, 3\,\% of the cells): the
        Bannerghatta hills, mostly ESRI rangeland;
  \item \textbf{the low western valley} (mean elevation \SI{759}{\m}, 3\,\% of the cells),
        where cropland and built land are mixed.
\end{itemize}

\begin{figure}[ht]
  \centering
  \includegraphics[width=\textwidth]{chart_k_scan}
  \caption{Choosing the number of clusters (2023 features). The dashed line marks $k=7$.}
  \label{fig:kscan}
\end{figure}

\paragraph{From clusters to classes.} Each cluster is labelled from its centroid
(Figure~\ref{fig:profiles}): water-dominated clusters are excluded; built-up where built land is
both above 30\,\% and the largest land group; forest where trees exceed 50\,\%, or where trees
and rangeland together exceed 50\,\% on slopes above \ang{8}; everything else is usable, with
cropland-dominated clusters flagged as agricultural. The second forest rule is what keeps the
Bannerghatta scrub forest out of the usable class. Finally, the exclusion mask (water, wetland,
slopes over \ang{15}, the national park) is applied cell by cell.

\begin{figure}[ht]
  \centering
  \includegraphics[width=0.95\textwidth]{chart_cluster_profiles}
  \caption{Cluster centroids (2023): mean shares, elevation and slope, coloured by z-score.
           Each row is named with the class its rule gave and its share of the cells.}
  \label{fig:profiles}
\end{figure}

\paragraph{The three-class map.} Figure~\ref{fig:class3} shows the result for 2023. Built-up land
covers \SI{348}{\km\squared}, usable land \SI{138}{\km\squared}, forest \SI{29}{\km\squared}, and
\SI{66}{\km\squared} are excluded. Between the 2018 and 2023 maps, built-up land grows by
\SI{42}{\km\squared} and usable land shrinks by \SI{44}{\km\squared}
(Figure~\ref{fig:class-areas}).

Against ESRI's own classes, collapsed to the same three, the map agrees on 96.6\,\% of the land
cells. This is expected, since the ESRI fractions are inputs, and it mainly shows that the
cluster rules don't distort them. Against the independent WorldCover map the agreement is
64.0\,\%. Almost all of the difference comes from the products themselves: WorldCover calls only
33\,\% of the area built-up, against 62\,\% in ESRI.

\begin{figure}[ht]
  \centering
  \includegraphics[width=0.7\textwidth]{map_class3_final}
  \caption{Three-class land map, 2023 (final run).}
  \label{fig:class3}
\end{figure}

\begin{figure}[ht]
  \centering
  \includegraphics[width=0.65\textwidth]{chart_class_areas}
  \caption{Class areas in the 2018 (validation) and 2023 (final) maps.}
  \label{fig:class-areas}
\end{figure}

\subsection{Building similarity}\label{sec:res-similarity}

\paragraph{Single-building queries.} Figure~\ref{fig:queries} shows two of the four example
queries from notebook~05, for a school and an industrial building. The most similar usable cells
lie at the edges of the built-up area, often several kilometres from A, wherever the mix of
built-up land, roads and terrain in the 250 and \SI{500}{\m} rings matches A's. Each result
comes with an explanation: for the industrial building's best match, the distances to built-up
land, water and a major road differ by less than 0.1 standard deviations, while the built-up
share within \SI{250}{\m} differs most (0.6 standard deviations: 80\,\% instead of 100\,\%). Land-cover shares that are zero at
both cells are listed last, since such matches say nothing. Euclidean and cosine distance share
14--18 of the top 20 cells; cosine ignores how extreme a profile is, so it is kept only as a
comparison.

\begin{figure}[ht]
  \centering
  \begin{subfigure}[t]{0.49\textwidth}
    \includegraphics[width=\textwidth]{map_similarity_query_school_1019}
  \end{subfigure}\hfill
  \begin{subfigure}[t]{0.49\textwidth}
    \includegraphics[width=\textwidth]{map_similarity_query_industrial_274244}
  \end{subfigure}
  \caption{Building similarity queries: the 20 usable cells whose surroundings are most similar
           to a school (left) and an industrial building (right), A marked with a star.}
  \label{fig:queries}
\end{figure}

\paragraph{Sanity checks.} For 200 random buildings, a cell within one cell of the building is in
the top 1\,\% of all cells even when the building's own cell is left out: surroundings vary
smoothly in space, so the profile identifies its location. Reference cells held out by
\SI{2}{\km} block are ranked above the other cells with AUC 0.895 by the nearest-neighbour
similarity, 0.984 by a random forest and 0.887 by distance to built-up alone. Keeping both ring
scales beats the \SI{500}{\m} rings alone (0.895 vs 0.871).

\paragraph{Aggregate similarity.} Averaged over the ten nearest reference profiles, the similarity
of every candidate cell reaches \textbf{AUC 0.773} on the growth by 2022/23
(Table~\ref{tab:auc}), above distance to built-up land, although it never sees a growth label.

\subsection{Learned models}\label{sec:res-rf}

The \textbf{change RF} reaches \textbf{AUC 0.779}, slightly above distance to built-up land
(Table~\ref{tab:auc}). Its most important input is the built-up share within \SI{500}{\m},
followed by trees and crops nearby and the distance to built-up land; terrain and water add
little (Figure~\ref{fig:importance}). Scored on its own training cells, it would claim AUC
0.997: most of its 2020/21 positives also grew by 2022/23, so in-sample scores would leak the
answer, and out-of-fold scoring is essential. The \textbf{state RF} separates built from
non-built cells almost perfectly (spatial-CV AUC 0.998), but that skill does not transfer to
predicting the next growth (0.702).

\begin{figure}[ht]
  \centering
  \includegraphics[width=0.85\textwidth]{chart_rf_importance}
  \caption{What the change RF relies on: permutation importance (AUC drop on held-out folds) and
           impurity importance.}
  \label{fig:importance}
\end{figure}

\subsection{Combined score, classes and candidate sites}\label{sec:res-suitability}

The weight chosen on the 2020/21 training label is 0.3 for the similarity and 0.7 for the change
RF. In the final run, the 13{,}814 usable cells fall into 1{,}382 S1, 2{,}762 S2, 4{,}144 S3 and
5{,}526 N cells (Figure~\ref{fig:suitability}). Highly suitable land concentrates on the southern
and eastern fringe of the built-up area, along the Hosur Road corridor, and in the low western
valley.

Merging edge-adjacent S1 cells gives 557 candidate patches covering 1{,}382\,ha; 253 are mostly
cropland and flagged as agricultural. Half of the patches are single \SI{1}{\hectare} cells,
often pockets inside built-up areas, so the top sites are taken from patches of at least
\SI{5}{\hectare}. The 20 top sites (Table~\ref{tab:sites}, 138\,ha) are all gently sloping
(\ang{1}--\ang{4}), directly next to existing built-up land, and covered by rangeland or crops.
None of them overlaps an OSM lake (0.06\,ha in total) or the national park.

\begin{figure}[ht]
  \centering
  \includegraphics[width=0.72\textwidth]{map_suitability_classes}
  \caption{Suitability classes in the final run (2023 data) with the outlines of the candidate
           patches. Forest, built-up and excluded land is always N.}
  \label{fig:suitability}
\end{figure}

\begin{table}[ht]
  \centering\small
  \caption{The ten best candidate sites (patches of at least \SI{5}{\hectare}, ranked by mean
           score). Distances are means over the patch's cells.}
  \label{tab:sites}
  \begin{tabular}{@{}rrrrrrll@{}}
    \toprule
    Rank & Area (ha) & Score & Slope (\textdegree) & Major road (m) & Built-up (m) & Cover & Location \\
    \midrule
    1 & 5 & 0.98 & 1.9 & 371 & 56 & rangeland & 77.64\,E, 12.86\,N \\
    2 & 6 & 0.98 & 1.6 & 1{,}621 & 55 & rangeland & 77.69\,E, 12.78\,N \\
    3 & 5 & 0.97 & 1.9 & 226 & 55 & crops & 77.73\,E, 12.86\,N \\
    4 & 5 & 0.97 & 3.9 & 572 & 57 & rangeland & 77.53\,E, 12.90\,N \\
    5 & 6 & 0.96 & 1.8 & 358 & 55 & rangeland & 77.71\,E, 12.76\,N \\
    6 & 6 & 0.96 & 1.6 & 161 & 55 & rangeland & 77.70\,E, 12.86\,N \\
    7 & 7 & 0.96 & 1.5 & 503 & 55 & rangeland & 77.66\,E, 12.76\,N \\
    8 & 6 & 0.96 & 1.5 & 1{,}265 & 55 & crops & 77.67\,E, 12.79\,N \\
    9 & 6 & 0.96 & 1.4 & 331 & 55 & rangeland & 77.71\,E, 12.74\,N \\
    10 & 5 & 0.95 & 1.3 & 412 & 61 & rangeland & 77.73\,E, 12.79\,N \\
    \bottomrule
  \end{tabular}
\end{table}

\subsection{Validation of all scores}\label{sec:res-validation}

Table~\ref{tab:auc} compares every score on the same 12{,}777 candidate cells, of which 1{,}125
were built in both 2022 and 2023. The \textbf{combined score} is the best: AUC 0.797, and its top
10\,\% of cells contain 3.45 times as much growth as a random tenth. Both PRD targets are met
(best model at least as good as distance to built-up; lift of at least 2 in the top 10\,\%). The
S1 and S2 classes of the validation run hold 65\,\% of the growth on 25\,\% of the evaluated
cells.

\begin{table}[ht]
  \centering\small
  \caption{Validation run: every score computed with 2018 data, tested on persistent growth by
           2022/23 (12{,}777 candidate cells, 1{,}125 grew). RF scores are out of fold.}
  \label{tab:auc}
  \begin{tabular}{@{}lccccc@{}}
    \toprule
    Score & AUC & Lift top 10\,\% & Lift top 30\,\% & AUC adjacent & AUC outlying \\
    \midrule
    Combined score & \textbf{0.797} & \textbf{3.45} & \textbf{2.40} & \textbf{0.822} & 0.512 \\
    Change RF & 0.779 & 3.23 & 2.27 & 0.803 & 0.512 \\
    Building similarity & 0.773 & 2.99 & 2.22 & 0.796 & 0.511 \\
    Distance to built-up (baseline) & 0.763 & 3.04 & 2.15 & 0.792 & 0.439 \\
    State RF & 0.702 & 3.29 & 1.92 & 0.722 & 0.484 \\
    MCDA / AHP & 0.636 & 1.86 & 1.64 & 0.655 & 0.420 \\
    MCDA without land cover & 0.614 & 1.73 & 1.46 & 0.630 & 0.430 \\
    Random & 0.499 & 1.00 & 1.02 & 0.501 & 0.477 \\
    \bottomrule
  \end{tabular}
\end{table}

\begin{figure}[ht]
  \centering
  \begin{subfigure}[t]{0.49\textwidth}
    \includegraphics[width=\textwidth]{chart_roc}
    \caption{ROC curves.}\label{fig:roc}
  \end{subfigure}\hfill
  \begin{subfigure}[t]{0.49\textwidth}
    \includegraphics[width=\textwidth]{chart_toc}
    \caption{Total operating characteristic.}\label{fig:toc}
  \end{subfigure}
  \caption{Predicting growth from 2018.}
\end{figure}

\paragraph{Growth type.} Every score finds growth next to existing built-up land well (AUC
0.79--0.82 for the best ones) but growth away from it only at chance level (0.51); distance to
built-up and MCDA even fall below chance there (Figure~\ref{fig:growth-type}). This matches the
finding that diffusional growth is much harder to predict than aggregated growth
\citep{wang2021}. Only 91 of the 1{,}125 grown cells are outlying, so this hardly moves the
overall AUC.

\paragraph{Ablation.} Dropping the built-up context (built-up shares in the rings and distance to
built-up) costs the most: $-0.034$ AUC for the change RF and $-0.057$ for the similarity
(Figure~\ref{fig:ablation}). Roads and water add nothing once the built-up context is known.

\begin{figure}[ht]
  \centering
  \begin{subfigure}[t]{0.52\textwidth}
    \includegraphics[width=\textwidth]{chart_growth_type}
    \caption{AUC by growth type.}\label{fig:growth-type}
  \end{subfigure}\hfill
  \begin{subfigure}[t]{0.46\textwidth}
    \includegraphics[width=\textwidth]{chart_ablation}
    \caption{AUC change when a feature group is dropped.}\label{fig:ablation}
  \end{subfigure}
  \caption{Where the models succeed and what they rely on.}
\end{figure}

\subsection{MCDA / AHP baseline}\label{sec:res-mcda}

The classic method uses five criteria, each rescaled to 0--1: distance to existing built-up land,
distance to a major road, slope, land cover (rangeland and bare ground easiest to develop, then
crops, then trees) and a setback from water. The AHP pairwise matrix gives built-up and road
distance 0.30 each, slope and land cover 0.16 each and water 0.09. Its consistency ratio of 0.003
is far below Saaty's 0.1 limit \citep{saaty1980}.

Tested against persistent growth from 2018 to 2022/23, the overlay reaches \textbf{AUC 0.636}
(Table~\ref{tab:auc}), with a lift of only 1.86 in the top 10\,\%, below a single criterion, distance to built-up land (0.763). Scored one at
a time, only that criterion carries real signal: slope reaches 0.612, and roads, land cover and
water stay between 0.52 and 0.55, close to chance. The AHP weights give the weak criteria 70\,\%
of the weight, which dilutes the one strong signal. We deliberately did not tune the weights on
the growth label, which is the test set. A fixed, defensible expert overlay is exactly what this
baseline stands for, and it illustrates the weakness the literature describes: the weights decide
the map, and nothing checks them against outcomes \citep{malczewski2004}.

\subsection{Sensitivity}\label{sec:res-sensitivity}

Three settings were varied (Figure~\ref{fig:sensitivity}):
\begin{itemize}
  \item \textbf{Slope threshold} (\ang{10}, \ang{15}, \ang{20}): the number of slope-excluded
        cells changes from 1{,}685 to 158, but every AUC moves by less than 0.01.
  \item \textbf{Number of clusters} ($k=5$ to 8): class shares change by about one percentage
        point and agreement with ESRI stays at 96.6--98.8\,\%. $k=5$ agrees slightly better
        with ESRI but loses the steep-vegetation cluster, so forest falls from 5.0\,\% to
        3.9\,\% of the area.
  \item \textbf{Cell size} (50, 100, \SI{200}{\m}; features and labels rebuilt each time): AUCs
        are not comparable across sizes, because the growth label itself changes with the cell,
        and at \SI{200}{\m} only 148 cells grew.
\end{itemize}
What holds in every setting is the ranking: \textbf{change RF $>$ distance to built-up $>$
MCDA}.

\begin{figure}[ht]
  \centering
  \includegraphics[width=0.95\textwidth]{chart_sensitivity}
  \caption{Temporal AUC by slope threshold and by cell size.}
  \label{fig:sensitivity}
\end{figure}

% Discussion and limitations (P8.7)
\section{Discussion and limitations}\label{sec:discussion}

\paragraph{Learned scores beat the expert overlay, but not by much against the simplest rule.}
The combined score improves on distance to built-up land by 0.034 AUC and on the classic AHP
overlay by 0.161. The first gap is modest because growth in Bengaluru South is overwhelmingly
edge expansion: 92\,\% of the grown cells touch land that was already built in 2018, and a single
distance already ranks such cells well. The value of the learned scores lies in the ranking
\emph{among} cells next to built-up land, where they add the composition of the neighbourhood
(how much built-up, cropland and scrub lies within 250 and \SI{500}{\m}) and terrain. Our AUC
is close to the 0.82 reported for a random forest tested on later change in Rasht
\citep{ahmadlou2016}, and far below the near-perfect values of studies that validate on a
random split of a single date \citep{li2024}: the difference comes from how the test is set
up, not from a better model.

\paragraph{Why the expert overlay does worse than one of its own criteria.} Only one of the five
MCDA criteria carries real signal on its own. The pairwise comparisons are internally
consistent (CR 0.003), but consistency says nothing about relevance: they put 70\,\% of the
weight on criteria that are close to random for this question. This is the subjectivity
\citet{malczewski2004} and \citet{mosadeghi2015} warn about, made measurable by a temporal test.

\paragraph{Similarity without labels works.} The building similarity never sees a growth label,
yet it predicts growth almost as well as the random forest trained on recent change (0.773 vs
0.779), and their combination is better than either. It also explains itself: each match
comes with the features where A and B agree and where they differ, which a planner can check.

\paragraph{Outlying growth is not predicted.} No score does better than chance for growth away
from existing built-up land, as \citet{wang2021} found for diffusional growth. Such growth
follows decisions the data do not record: land deals, new industrial estates, approvals.

\paragraph{Limitations.}
\begin{itemize}
  \item \textbf{Land-cover data.} ESRI and WorldCover disagree on built-up land (62\,\% vs
        33\,\% of the area): ESRI generalises urban areas, and it calls part of the
        Bannerghatta forest rangeland \citep{venter2022}. Growth labels inherit its errors; the
        two-year rule removes flicker but not systematic bias.
  \item \textbf{OpenStreetMap.} Road mapping more than doubled between the snapshots, mostly
        through mapping effort, so each run uses its own snapshot. OSM buildings were sparse in
        the south in 2018 \citep{herfort2023} and almost none are typed, so similarity per
        building type was not possible.
  \item \textbf{Resolution.} The 100\,m cells mix land uses; half of the top-class patches are
        single cells inside built-up areas. Cell size changes the growth label itself, so
        results at different resolutions are not directly comparable.
  \item \textbf{Past growth is not good planning.} Every score learns where building
        \emph{has} happened. Water, steep slopes and the national park are excluded, but
        zoning, flood risk, infrastructure capacity and ownership are not in the data. The
        suitability map shows where growth is \emph{likely}, as an input to planning, not a
        recommendation.
\end{itemize}

% Conclusion and future work (P8.7)
\section{Conclusion and future work}\label{sec:conclusion}

This project built a land suitability analysis for urban development in Bengaluru South
entirely in Python, from open land-cover, terrain and OpenStreetMap data, and tested every
method against the growth that actually happened. It answers the three questions it set out
with:

\begin{enumerate}
  \item \textbf{Land classes.} Clustering 100\,m cells on land-cover shares and terrain gives
        seven interpretable clusters, mapped to built-up (348\,km\textsuperscript{2} in 2023),
        forest (29\,km\textsuperscript{2}, including the steep scrub forest of Bannerghatta) and
        usable land (138\,km\textsuperscript{2}). The map agrees with the ESRI classes on
        96.6\,\% of land cells.
  \item \textbf{Building similarity.} For any building, the usable locations whose surroundings
        look most like its own can be found and explained feature by feature. Averaged over
        reference buildings, this similarity predicts growth from 2018 to 2022/23 with
        AUC 0.773, without ever seeing a growth label.
  \item \textbf{Suitability.} Combining the similarity with a random forest trained on recent
        change gives the best score (AUC 0.797; the top 10\,\% of cells hold 3.45 times their
        share of growth), clearly ahead of a classic AHP overlay (0.636) and ahead of the
        distance to built-up land (0.763). The final map grades usable land into classes
        S1--S3 and lists the 20 best candidate sites.
\end{enumerate}

The main lesson is about validation: strict rules against leakage (no own-cell land cover, data
frozen at the baseline date, out-of-fold scores) turn near-perfect accuracies into honest ones,
and simple baselines show how much of the signal is just ``next to existing built-up land''.

\paragraph{Future work.}
\begin{itemize}
  \item Add drivers of growth away from the built-up edge: planned layouts, industrial estates,
        land prices, distance to the city centre and to employment.
  \item Re-run the analysis with newer years as they become available, moving the baseline
        forward to test whether the models generalise in time.
  \item Use building footprints from other sources (e.g. Open Buildings) to type buildings, so
        that similarity can be computed per building type.
  \item Test finer grids with a growth label defined independently of the cell size.
\end{itemize}


\bibliographystyle{plainnat}
\bibliography{references}

\end{document}
